Wyniki przedstawiamy od najogólniejszego pytania, czyli który sposób losowania daje najlepszy zespół (H3). Następnie opisujemy mechanizm, czyli wpływ rozmiaru i~pokrywania się próbek na siłę i~niezależność drzew, a~potem hipotezy o~miarach próbek (H1, H2). Na końcu omawiamy eksperyment dodatkowy i~porównanie lasu losowego z~baggingiem. Średnie i~odchylenia standardowe metryki głównej z~50 podziałów dla każdego zbioru i~wariantu zawiera Dodatek~\ref{app:szczegoly}.

% --------------------------------------------------------------------
\subsection{Porównanie wariantów z~wariantem odniesienia (H3)}
\label{sec:wyniki-h3}

Test Friedmana odrzuca hipotezę o~równoważności 15 wariantów w~obu zadaniach (Tabela~\ref{tab:friedman}), więc sposób losowania wpływa na jakość zespołu. Średnie rangi (Rysunek~\ref{fig:rangi}) pokazują jednak, że wariant odniesienia należy do najlepszych. W~regresji ma najlepszą średnią rangę (4.3). W~klasyfikacji lepsze rangi mają tylko losowanie warstwowe S4 (3.9) i~losowanie bez zwracania z~$s=0.63$ (4.0), czyli warianty najbardziej do niego podobne.

\begin{table}[!htbp]
\centering\small
\caption{Test Friedmana dla 15 wariantów (las losowy) i~jego wersja Imana--Davenporta.}
\label{tab:friedman}
\begin{tabular}{lrrrrr}
\toprule
\textbf{Zadanie} & \textbf{zbiory} & $\chi^2_F$ & $p$ & $F_{ID}$ & $p$ \\
\midrule
Klasyfikacja & 12 & 75.5 & <0.001 & 8.99 & <0.001 \\
Regresja & 12 & 76.9 & <0.001 & 9.29 & <0.001 \\
\bottomrule
\end{tabular}
\end{table}


\begin{figure}[!htbp]
\centering
\includegraphics[width=\textwidth]{05_average_ranks.png}
\caption{Średnie rangi 15 wariantów w~teście Friedmana (12 zbiorów w~każdym zadaniu, las losowy). Czarną ramką oznaczono wariant odniesienia S1~$r=1$.}
\label{fig:rangi}
\end{figure}

Porównania z~wariantem odniesienia zawiera Tabela~\ref{tab:kontrola}.

Żaden wariant nie jest istotnie lepszy od wariantu odniesienia, ani w~klasyfikacji, ani w~regresji, ani w~analizie łącznej 24 zbiorów. Dwa warianty o~najlepszych rangach w~klasyfikacji (S4 i~S2~$s=0.63$) poprawiają macro-F1 o~medianę $0.09$ i~$0.04$\pp{}. Test bayesowski przypisuje obu prawdopodobieństwo $1.00$ praktycznej równoważności z~wariantem odniesienia, więc wynik wskazuje na brak różnicy praktycznej, a~nie tylko na brak istotności.

Istotnie gorsze są warianty z~dużym wspólnym rdzeniem: S3 z~$\gamma\ge0.25$ w~regresji i~z $\gamma=0.75$ w~klasyfikacji. Przegrywają one w~11 lub 12 z~12 zbiorów, a~w~regresji mediana wzrostu RMSE wynosi od 2.8 do 8.6\%.

Małe próbki (S1 z~$r\le0.4$, S2 z~$s\le0.33$) przegrywają z~wariantem odniesienia w~9--11 z~12 zbiorów, a~w~regresji mediana straty sięga 10\% RMSE. Przy 12 zbiorach wynik nie jest jednak istotny po korekcie Holma ($p_H$ od 0.08 do 0.16). Zgadza się to z~analizą mocy (Tabela~\ref{tab:moc}): wystarczy kilka zbiorów, w~których mniejsze próbki pomagają, aby test nie odrzucił hipotezy zerowej. W~analizie łącznej 24 zbiorów te same warianty są istotnie gorsze ($p_H\le0.014$), a~test bayesowski daje prawdopodobieństwo od 0.93 do 1.00, że są gorsze o~więcej niż szerokość ROPE.

Strata przy najmniejszych próbkach przekracza odchylenie standardowe wyniku wariantu odniesienia między podziałami w~5 zbiorach klasyfikacyjnych i~w~8 regresyjnych. Dla wariantów najbliższych wariantowi odniesienia (S1~$r=0.8$, S2~$s\ge0.55$, S4) różnica nie przekracza tego odchylenia w~żadnym zbiorze.

Mediana stosunku odchyleń standardowych $\sigma/\sigma_0$ wynosi od 0.99 do 1.09 dla wszystkich wariantów z~wyjątkiem S3~$\gamma=0.75$ w~regresji (1.20). Żaden wariant nie jest więc wyraźnie stabilniejszy od wariantu odniesienia, a~duży wspólny rdzeń, oprócz gorszej średniej, zwiększa zmienność wyniku między podziałami.

{\scriptsize\setlength{\tabcolsep}{4pt}
\begin{longtable}{lrrrrrrcrr}
\caption{Porównanie wariantów z~wariantem odniesienia S1~$r=1$ (las losowy). Ranga to średnia ranga w~teście Friedmana (1 oznacza najlepszy z~15 wariantów). Mediana to mediana wyniku względnego po zbiorach; wartość dodatnia oznacza wariant lepszy od wariantu odniesienia. W/P to liczba zbiorów, w~których wariant wygrał i~przegrał. $r_{rb}$ to dwuseryjny współczynnik korelacji rang, a~$p_H$ to wartość $p$ po korekcie Holma (pogrubiono $p_H<0.05$). Kolumna Bayes zawiera prawdopodobieństwa a~posteriori, że wariant jest gorszy, praktycznie równoważny (ROPE $\pm1$) i~lepszy. Kolumna ,,Istotne prakt.'' podaje liczbę zbiorów, w~których różnica na korzyść i~na niekorzyść wariantu przekracza odchylenie standardowe wyniku wariantu odniesienia między 50 podziałami. $\sigma/\sigma_0$ to mediana po zbiorach stosunku odchyleń standardowych metryki głównej wariantu i~wariantu odniesienia między podziałami; wartość poniżej 1 oznacza wariant stabilniejszy.}
\label{tab:kontrola}\\
\toprule
\textbf{Wariant} & \textbf{Ranga} & \textbf{Mediana} & \textbf{W/P} & $r_{rb}$ & $p$ & $p_H$ & \textbf{Bayes} & \textbf{Istotne prakt.} & $\sigma/\sigma_0$ \\
\endfirsthead
\toprule
\textbf{Wariant} & \textbf{Ranga} & \textbf{Mediana} & \textbf{W/P} & $r_{rb}$ & $p$ & $p_H$ & \textbf{Bayes} & \textbf{Istotne prakt.} & $\sigma/\sigma_0$ \\
\endhead
\midrule
\multicolumn{10}{l}{\emph{Klasyfikacja} -- $\Delta$ macro-F1 [p.p.]}\\
S1~$r=0.2$ & 12.2 & -1.62 & 1/11 & -0.85 & 0.007 & 0.089 & 0.88 / 0.12 / 0.00 & 1/5 & 1.07 \\
S1~$r=0.4$ & 10.6 & -0.58 & 1/11 & -0.77 & 0.016 & 0.161 & 0.23 / 0.77 / 0.00 & 1/1 & 1.05 \\
S1~$r=0.6$ & 8.0 & -0.28 & 3/9 & -0.67 & 0.042 & 0.382 & 0.00 / 1.00 / 0.00 & 1/0 & 1.04 \\
S1~$r=0.8$ & 5.8 & -0.13 & 5/7 & -0.08 & 0.850 & 1.000 & 0.00 / 1.00 / 0.00 & 0/0 & 1.04 \\
S2~$s=0.18$ & 13.0 & -1.58 & 1/11 & -0.85 & 0.007 & 0.089 & 0.90 / 0.10 / 0.00 & 1/5 & 1.01 \\
S2~$s=0.33$ & 11.4 & -0.75 & 1/11 & -0.82 & 0.009 & 0.102 & 0.27 / 0.73 / 0.00 & 1/1 & 1.01 \\
S2~$s=0.45$ & 6.2 & -0.13 & 5/7 & -0.38 & 0.266 & 1.000 & 0.00 / 1.00 / 0.00 & 1/0 & 1.01 \\
S2~$s=0.55$ & 5.3 & -0.03 & 6/6 & -0.18 & 0.622 & 1.000 & 0.00 / 1.00 / 0.00 & 0/0 & 0.99 \\
S2~$s=0.63$ & 4.0 & +0.04 & 8/4 & +0.36 & 0.301 & 1.000 & 0.00 / 1.00 / 0.00 & 0/0 & 1.03 \\
S3~$\gamma=0$ & 7.8 & -0.23 & 3/9 & -0.62 & 0.064 & 0.512 & 0.00 / 1.00 / 0.00 & 0/0 & 1.06 \\
S3~$\gamma=0.25$ & 7.1 & -0.12 & 4/8 & -0.38 & 0.266 & 1.000 & 0.00 / 1.00 / 0.00 & 0/0 & 1.08 \\
S3~$\gamma=0.5$ & 8.7 & -0.19 & 3/9 & -0.62 & 0.064 & 0.512 & 0.00 / 1.00 / 0.00 & 0/1 & 1.04 \\
S3~$\gamma=0.75$ & 10.8 & \textbf{-0.48} & 1/11 & -0.97 & <0.001 & \textbf{0.014} & 0.18 / 0.82 / 0.00 & 0/1 & 1.09 \\
S4 & 3.9 & +0.09 & 8/4 & +0.62 & 0.064 & 0.512 & 0.00 / 1.00 / 0.00 & 0/0 & 1.03 \\
\midrule
\multicolumn{10}{l}{\emph{Regresja} -- poprawa RMSE [\%]}\\
S1~$r=0.2$ & 11.3 & -10.16 & 3/9 & -0.85 & 0.007 & 0.075 & 1.00 / 0.00 / 0.00 & 0/8 & 1.01 \\
S1~$r=0.4$ & 8.7 & -4.40 & 3/9 & -0.79 & 0.012 & 0.110 & 0.99 / 0.01 / 0.00 & 0/7 & 0.99 \\
S1~$r=0.6$ & 6.3 & -1.82 & 4/8 & -0.69 & 0.034 & 0.205 & 0.91 / 0.08 / 0.01 & 0/2 & 0.99 \\
S1~$r=0.8$ & 4.6 & -0.44 & 5/7 & -0.46 & 0.176 & 0.353 & 0.11 / 0.89 / 0.00 & 0/0 & 0.99 \\
S2~$s=0.18$ & 11.3 & -10.20 & 3/9 & -0.85 & 0.007 & 0.075 & 1.00 / 0.00 / 0.00 & 0/8 & 1.01 \\
S2~$s=0.33$ & 9.0 & -4.45 & 3/9 & -0.74 & 0.021 & 0.168 & 0.99 / 0.01 / 0.00 & 0/7 & 1.00 \\
S2~$s=0.45$ & 6.7 & -1.99 & 4/8 & -0.72 & 0.027 & 0.188 & 0.90 / 0.10 / 0.00 & 0/2 & 1.00 \\
S2~$s=0.55$ & 5.9 & -0.75 & 4/8 & -0.67 & 0.042 & 0.205 & 0.36 / 0.64 / 0.00 & 0/0 & 1.01 \\
S2~$s=0.63$ & 5.9 & -0.60 & 1/11 & -0.69 & 0.034 & 0.205 & 0.04 / 0.96 / 0.00 & 0/0 & 1.01 \\
S3~$\gamma=0$ & 6.2 & -1.07 & 4/8 & -0.67 & 0.042 & 0.205 & 0.57 / 0.43 / 0.00 & 0/2 & 1.00 \\
S3~$\gamma=0.25$ & 9.8 & \textbf{-2.77} & 1/11 & -0.97 & <0.001 & \textbf{0.012} & 0.98 / 0.02 / 0.00 & 0/2 & 1.02 \\
S3~$\gamma=0.5$ & 11.7 & \textbf{-4.54} & 0/12 & -1.00 & <0.001 & \textbf{0.007} & 1.00 / 0.00 / 0.00 & 0/7 & 1.06 \\
S3~$\gamma=0.75$ & 13.8 & \textbf{-8.64} & 0/12 & -1.00 & <0.001 & \textbf{0.007} & 1.00 / 0.00 / 0.00 & 0/7 & 1.20 \\
S4 & 4.6 & -0.03 & 4/8 & -0.23 & 0.519 & 0.519 & 0.00 / 1.00 / 0.00 & 0/0 & 1.01 \\
\midrule
\multicolumn{10}{l}{\emph{Łącznie (24 zbiory)} -- wynik względny [p.p. / \%]}\\
S1~$r=0.2$ & -- & \textbf{-3.74} & 4/20 & -0.79 & <0.001 & \textbf{0.003} & 1.00 / 0.00 / 0.00 & -- & -- \\
S1~$r=0.4$ & -- & \textbf{-1.05} & 4/20 & -0.70 & 0.002 & \textbf{0.014} & 0.93 / 0.07 / 0.00 & -- & -- \\
S1~$r=0.6$ & -- & -0.41 & 7/17 & -0.57 & 0.013 & 0.088 & 0.28 / 0.72 / 0.00 & -- & -- \\
S1~$r=0.8$ & -- & -0.17 & 10/14 & -0.25 & 0.290 & 0.579 & 0.00 / 1.00 / 0.00 & -- & -- \\
S2~$s=0.18$ & -- & \textbf{-3.85} & 4/20 & -0.79 & <0.001 & \textbf{0.003} & 1.00 / 0.00 / 0.00 & -- & -- \\
S2~$s=0.33$ & -- & \textbf{-1.18} & 4/20 & -0.71 & 0.001 & \textbf{0.013} & 0.94 / 0.06 / 0.00 & -- & -- \\
S2~$s=0.45$ & -- & -0.36 & 9/15 & -0.50 & 0.031 & 0.157 & 0.28 / 0.72 / 0.00 & -- & -- \\
S2~$s=0.55$ & -- & -0.11 & 10/14 & -0.47 & 0.042 & 0.170 & 0.02 / 0.98 / 0.00 & -- & -- \\
S2~$s=0.63$ & -- & -0.27 & 9/15 & -0.43 & 0.069 & 0.207 & 0.00 / 1.00 / 0.00 & -- & -- \\
S3~$\gamma=0$ & -- & -0.27 & 7/17 & -0.55 & 0.016 & 0.098 & 0.05 / 0.95 / 0.00 & -- & -- \\
S3~$\gamma=0.25$ & -- & \textbf{-0.52} & 5/19 & -0.78 & <0.001 & \textbf{0.004} & 0.56 / 0.44 / 0.00 & -- & -- \\
S3~$\gamma=0.5$ & -- & \textbf{-1.47} & 3/21 & -0.90 & <0.001 & \textbf{<0.001} & 0.92 / 0.08 / 0.00 & -- & -- \\
S3~$\gamma=0.75$ & -- & \textbf{-3.81} & 1/23 & -0.99 & <0.001 & \textbf{<0.001} & 1.00 / 0.00 / 0.00 & -- & -- \\
S4 & -- & +0.02 & 12/12 & +0.21 & 0.390 & 0.579 & 0.00 / 1.00 / 0.00 & -- & -- \\
\bottomrule
\end{longtable}
}


Rysunek~\ref{fig:zbiory} pokazuje wynik każdego wariantu w~każdym zbiorze. Kierunek wpływu rozmiaru próbki zależy od zbioru. Mniejsze próbki poprawiają wynik w~zbiorach z~dużym szumem w~zmiennej docelowej: Adult (do $+1.5$\pp{} macro-F1), Abalone ($+2.4\%$), Diabetes ($+3.2\%$) i~Wine Quality ($+0.6\%$). Znacznie pogarszają go natomiast w~zbiorach regresyjnych o~niemal deterministycznej zależności i~niewielu cechach. W~zbiorach Airfoil, Energy Efficiency, Yacht i~Concrete RMSE przy $r=0.2$ rośnie o~25--48\%. W~klasyfikacji największą stratę ma Glass ($-16.7$\pp), najmniejszy zbiór, w~którym jest sześć klas, a~najrzadsza z~nich ma w~części treningowej około 7 obserwacji. Próbka z~$r=0.2$ często w~ogóle jej nie zawiera. W~tym samym zbiorze losowanie warstwowe daje największą poprawę ($+1.3$\pp), bo zapewnia obecność każdej klasy w~każdej próbce.

\begin{figure}[!htbp]
\centering
\includegraphics[width=\textwidth]{08_per_dataset.png}
\caption{Wynik każdego wariantu względem wariantu odniesienia w~każdym zbiorze (średnia z~50 podziałów, las losowy). Lewy panel: różnica macro-F1 w~punktach procentowych. Prawy panel: względna poprawa RMSE w~procentach. Kolor niebieski oznacza wariant lepszy od wariantu odniesienia.}
\label{fig:zbiory}
\end{figure}

% --------------------------------------------------------------------
\subsection{Rozmiar próbki i~pokrywanie się próbek}
\label{sec:wyniki-rozmiar}

Rysunek~\ref{fig:rozmiar} pokazuje dwa czynniki, które warianty zmieniają celowo. Lewy panel przedstawia wynik S1 i~S2 w~funkcji udziału różnych obserwacji $U$. Krzywe obu rodzin prawie się pokrywają, więc przy tej samej liczbie różnych obserwacji nie ma znaczenia, czy próbka zawiera powtórzenia. Potwierdzają to mediany w~analizie łącznej (Tabela~\ref{tab:kontrola}). S1~$r=0.2$ i~S2~$s=0.18$ tracą odpowiednio 3.74 i~3.85 jednostki, a~S1~$r=0.4$ i~S2~$s=0.33$ odpowiednio 1.05 i~1.18 jednostki. Mediana wyniku rośnie monotonicznie z~$U$, choć, jak pokazano wyżej, nie w~każdym zbiorze.

\begin{figure}[!htbp]
\centering
\includegraphics[width=\textwidth]{02_size_and_core.png}
\caption{Lewy panel: wynik względem wariantu odniesienia w~funkcji udziału różnych obserwacji $U$ dla S1 (losowanie ze zwracaniem) i~S2 (bez zwracania). Prawy panel: wynik S3 w~funkcji udziału wspólnego rdzenia $\gamma$ przy stałym rozmiarze $m=0.5n$. Cienkie linie odpowiadają poszczególnym zbiorom, grube medianie po zbiorach.}
\label{fig:rozmiar}
\end{figure}

W~prawym panelu Rysunku~\ref{fig:rozmiar} rozmiar próbki jest stały ($m=0.5n$), a~zmienia się tylko stopień pokrywania się próbek różnych drzew. W~regresji jakość zespołu maleje monotonicznie z~$\gamma$ w~każdym zbiorze. W~klasyfikacji wyraźny spadek występuje dopiero przy $\gamma=0.75$. Tabela~\ref{tab:s3} zawiera bezpośredni test drugiej części H2. Próbki niezależne ($\gamma=0$) dają lepszy zespół niż próbki w~75\% wspólne w~9 z~12 zbiorów klasyfikacyjnych ($p=0.007$) i~we wszystkich 12 regresyjnych ($p<0.001$, mediana $+7.4\%$ RMSE). W~klasyfikacji różnica jest jednak mała: według testu bayesowskiego z~prawdopodobieństwem 0.94 mieści się w~ROPE, czyli jest mniejsza niż 1\pp{} macro-F1. Sekcja~\ref{sec:rf-bagging} pokazuje, że wynika to z~losowania cech w~lesie losowym.

\begin{table}[!htbp]
\centering\small
\caption{Bezpośredni test drugiej części H2. Wariant S3 z~$\gamma=0$ (próbki niezależne) porównano z~$\gamma=0.75$ (75\% każdej próbki wspólne dla wszystkich drzew) przy tym samym rozmiarze $m=0.5n$. Dodatnia różnica oznacza, że mniejsze pokrywanie się próbek daje lepszy zespół. Kolumna Bayes zawiera kolejno prawdopodobieństwa a~posteriori, że lepsze jest $\gamma=0.75$, że różnica mieści się w~ROPE i~że lepsze jest $\gamma=0$.}
\label{tab:s3}
\begin{tabular}{lrrrrc}
\toprule
\textbf{Zadanie} & \textbf{zbiory} & \textbf{mediana różnicy} & \textbf{wygrane $\gamma=0$} & $p$ & \textbf{Bayes} \\
\midrule
Klasyfikacja & 12 & +0.31 & 9 & 0.007 & 0.00 / 0.94 / 0.06 \\
Regresja & 12 & +7.38 & 12 & <0.001 & 0.00 / 0.00 / 1.00 \\
Łącznie (24 zbiory) & 24 & +4.20 & 21 & <0.001 & 0.00 / 0.00 / 1.00 \\
\bottomrule
\end{tabular}
\end{table}


% --------------------------------------------------------------------
\subsection{Siła i~niezależność drzew}
\label{sec:wyniki-sila}

Rysunek~\ref{fig:sila} pokazuje położenie zespołów na płaszczyźnie wyznaczonej przez siłę drzew i~różnorodność ich predykcji. Obie wielkości standaryzowaliśmy w~obrębie zbioru, a~jeden punkt odpowiada jednemu wariantowi w~jednym zbiorze. Warianty S1 i~S2 leżą na jednej malejącej krzywej: mniejsza próbka daje drzewa bardziej różne, ale słabsze. Warianty S3 układają się w~poziomy pas, ponieważ przy prawie stałej sile drzew zmienia się tylko ich różnorodność, a~zespoły o~mniejszej różnorodności (większe $\gamma$) są gorsze. Wariant S3 rozdziela więc, zgodnie z~założeniem, niezależność drzew od ich siły.

\begin{figure}[!htbp]
\centering
\includegraphics[width=\textwidth]{03_strength_vs_diversity.png}
\caption{Średnia jakość drzew a~różnorodność ich predykcji (wartości standaryzowane w~obrębie zbioru), las losowy. Każdy punkt odpowiada jednemu wariantowi w~jednym zbiorze, a~kolor wypełnienia wynikowi zespołu względem wariantu odniesienia.}
\label{fig:sila}
\end{figure}

% --------------------------------------------------------------------
\subsection{H1: miary próbek a~różnice między drzewami jednego zespołu}
\label{sec:wyniki-h1}

Tabela~\ref{tab:h1} i~Rysunek~\ref{fig:h1} zawierają korelacje między miarami próbek a~własnościami drzew w~obrębie jednego zespołu wariantu odniesienia.

Pary drzew uczonych na bardziej różnych próbkach ($D_J$, $D_C$) dają bardziej różne predykcje. Korelacja jest dodatnia w~10--12 z~12 zbiorów i~istotna po korekcie Holma w~obu zadaniach i~w~obu konfiguracjach zespołu. Jest jednak bardzo słaba. Mediana $\rho$ wynosi 0.02--0.03, więc miara wyjaśnia rzędu 0.1\% zmienności różnorodności par. Różnica rozkładu zmiennej docelowej $D_Y$ wiąże się z~różnorodnością predykcji jeszcze słabiej i~istotnie tylko w~regresji ($\tilde\rho=0.007$).

Żadna z~miar reprezentatywności ($U$, $R_Y$, $R_X$) nie jest istotnie związana z~jakością pojedynczego drzewa. Dla $U$ korelacja jest dodatnia w~9 lub 10 zbiorach ($\tilde\rho\approx0.02$), ale nie jest istotna po korekcie Holma. Dla $R_Y$ i~$R_X$ znak korelacji zmienia się między zbiorami.

\begin{table}[!htbp]
\centering\small
\caption{H1: korelacje Spearmana wewnątrz zespołu w~wariancie odniesienia S1~$r=1$. Podano medianę po zbiorach (z~median po 50 podziałach), liczbę zbiorów z~dodatnią korelacją oraz $p_H$, czyli wartość $p$ testu Wilcoxona po zbiorach z~korektą Holma (6 testów w~zadaniu). W~regresji las losowy jest baggingiem, więc kolumny dla baggingu pominięto.}
\label{tab:h1}
\begin{tabular}{lrrrrrr}
\toprule
& \multicolumn{3}{c}{\textbf{Las losowy}} & \multicolumn{3}{c}{\textbf{Bagging}} \\
\cmidrule(lr){2-4}\cmidrule(lr){5-7}
\textbf{Korelacja} & $\tilde\rho$ & $\rho>0$ & $p_H$ & $\tilde\rho$ & $\rho>0$ & $p_H$ \\
\midrule
\multicolumn{7}{l}{\emph{Klasyfikacja}}\\
pary: $D_J$ ~ różn.\ predykcji & +0.018 & 10/12 & \textbf{0.015} & +0.033 & 12/12 & \textbf{0.003} \\
pary: $D_C$ ~ różn.\ predykcji & +0.018 & 10/12 & \textbf{0.015} & +0.031 & 12/12 & \textbf{0.003} \\
pary: $D_Y$ ~ różn.\ predykcji & +0.003 & 9/12 & 0.904 & +0.006 & 7/12 & 0.529 \\
drzewa: $U$ ~ jakość drzewa & +0.017 & 9/12 & 0.107 & +0.017 & 9/12 & 0.369 \\
drzewa: $R_Y$ ~ jakość drzewa & +0.004 & 8/12 & 0.904 & -0.011 & 4/12 & 0.529 \\
drzewa: $R_X$ ~ jakość drzewa & -0.016 & 4/12 & 0.904 & -0.004 & 5/12 & 0.529 \\
\midrule
\multicolumn{7}{l}{\emph{Regresja}}\\
pary: $D_J$ ~ różn.\ predykcji & +0.033 & 12/12 & \textbf{0.003} & \multicolumn{3}{c}{= las losowy} \\
pary: $D_C$ ~ różn.\ predykcji & +0.031 & 12/12 & \textbf{0.003} & \multicolumn{3}{c}{= las losowy} \\
pary: $D_Y$ ~ różn.\ predykcji & +0.007 & 10/12 & \textbf{0.037} & \multicolumn{3}{c}{= las losowy} \\
drzewa: $U$ ~ jakość drzewa & +0.023 & 10/12 & 0.127 & \multicolumn{3}{c}{= las losowy} \\
drzewa: $R_Y$ ~ jakość drzewa & -0.003 & 6/12 & 0.761 & \multicolumn{3}{c}{= las losowy} \\
drzewa: $R_X$ ~ jakość drzewa & -0.002 & 5/12 & 0.910 & \multicolumn{3}{c}{= las losowy} \\
\bottomrule
\end{tabular}
\end{table}


\begin{figure}[!htbp]
\centering
\includegraphics[width=\textwidth]{04_h1_within_ensemble.png}
\caption{H1 dla lasu losowego w~wariancie odniesienia: mediana (po 50 podziałach) korelacji Spearmana w~każdym zbiorze. Kolor niebieski: miary par próbek a~różnorodność predykcji par drzew. Kolor pomarańczowy: miary pojedynczej próbki a~jakość drzewa.}
\label{fig:h1}
\end{figure}

Przyczynę słabości tych związków widać na Rysunku~\ref{fig:pary}. W~jednym zespole odległość Jaccarda wszystkich 4950 par próbek mieści się w~wąskim przedziale od około 0.46 do 0.62, zgodnie z~koncentracją opisaną w~Sekcji~\ref{sec:koncentracja}. Różnorodność predykcji tych samych par zmienia się natomiast kilkukrotnie. Różnice między drzewami zespołu wynikają więc przede wszystkim z~tego, \emph{które} obserwacje (a~w~lesie losowym także które cechy) trafiły do próbki, a~nie z~tego, \emph{ile} obserwacji dwie próbki mają wspólnych.

\begin{figure}[!htbp]
\centering
\includegraphics[width=\textwidth]{09_pair_example.png}
\caption{Przykład H1 w~jednym podziale: odległość Jaccarda i~różnorodność predykcji dla 4950 par drzew w~zbiorach Breast Cancer (klasyfikacja, niezgodność) i~Diabetes (regresja, $1-$korelacja błędów).}
\label{fig:pary}
\end{figure}

% --------------------------------------------------------------------
\subsection{H2: miary próbek a~jakość zespołu}
\label{sec:wyniki-h2}

Tabela~\ref{tab:h2} zawiera korelacje między średnią wartością miary a~jakością zespołu po 15 wariantach, liczone w~każdym zbiorze osobno. Rysunek~\ref{fig:h2} pokazuje te korelacje dla poszczególnych zbiorów.

\textbf{Korelacja prosta} ma w~większości zbiorów ten sam znak. Wyższe $U$ idzie w~parze z~lepszym zespołem, a~większa różnorodność par ($D_J$, $D_C$, $D_Y$) i~gorsza reprezentatywność ($R_Y$, $R_X$) z~gorszym. Wszystkie te miary zmieniają się bowiem razem z~rozmiarem próbki. Kierunek nie jest jednak wspólny dla wszystkich zbiorów. W~Adult, Abalone, Diabetes i~Wine Quality znaki są odwrócone, a~są to te same zbiory, w~których mniejsze próbki poprawiają wynik (Rysunek~\ref{fig:zbiory}). Dlatego żadna korelacja prosta nie jest istotna po korekcie Holma, choć np.\ dla $U$ znak jest ten sam w~9 lub 10 z~12 zbiorów.

\textbf{Korelacja cząstkowa}, czyli korelacja po usunięciu związku ze średnią jakością drzew, jest dodatnia prawie we wszystkich zbiorach (dolny wiersz Rysunku~\ref{fig:h2}). W~regresji wszystkie miary różnorodności par są dodatnio związane z~jakością zespołu we wszystkich 12 zbiorach ($\tilde\rho$ od 0.58 do 0.69, $p_H=0.003$). Najsilniejszy związek ma różnorodność predykcji ($\tilde\rho=0.77$). Przy tej samej sile drzew bardziej różne próbki dają więc lepszy zespół, zgodnie z~oszacowaniem Breimana. Dodatnie korelacje cząstkowe $R_Y$ i~$R_X$ mają tę samą przyczynę. Przy ustalonej sile drzew mniej reprezentatywna próbka jest w~praktyce próbką mniejszą, a~więc bardziej różną od pozostałych. W~lesie losowym w~klasyfikacji korelacje cząstkowe są dodatnie w~9 lub 10 z~12 zbiorów, ale słabsze i~nieistotne po korekcie ($p_H\ge0.19$). W~baggingu (Tabela~\ref{tab:h2bag}) ten sam związek jest istotny i~występuje we wszystkich 12 zbiorach ($\tilde\rho$ od 0.54 do 0.62 dla miar par, $p_H=0.003$).

\begin{table}[!htbp]
\centering\small
\caption{H2: korelacja Spearmana między średnią wartością miary a~wynikiem zespołu po 15 wariantach, liczona w~każdym zbiorze osobno. Podano medianę po zbiorach, liczbę zbiorów z~dodatnią korelacją i~wartość $p_H$ testu Wilcoxona z~korektą Holma (7 miar). Korelacja cząstkowa uwzględnia średnią jakość drzew.}
\label{tab:h2}
\begin{tabular}{lrrrrrr}
\toprule
& \multicolumn{3}{c}{\textbf{korelacja}} & \multicolumn{3}{c}{\textbf{cząstkowa}} \\
\cmidrule(lr){2-4}\cmidrule(lr){5-7}
\textbf{Miara} & $\tilde\rho$ & $\rho>0$ & $p_H$ & $\tilde\rho$ & $\rho>0$ & $p_H$ \\
\midrule
\multicolumn{7}{l}{\emph{Klasyfikacja}}\\
$D_J$ & -0.56 & 2/12 & 0.196 & +0.28 & 9/12 & 0.205 \\
$D_C$ & -0.45 & 2/12 & 0.196 & +0.22 & 10/12 & 0.212 \\
$D_Y$ & -0.60 & 2/12 & 0.196 & +0.27 & 9/12 & 0.330 \\
$U$ & +0.83 & 10/12 & 0.188 & +0.22 & 9/12 & 0.212 \\
$R_Y$ & -0.75 & 2/12 & 0.196 & +0.08 & 7/12 & 1.000 \\
$R_X$ & -0.69 & 2/12 & 0.196 & +0.04 & 7/12 & 1.000 \\
różnorodność predykcji & -0.64 & 2/12 & 0.196 & +0.41 & 10/12 & 0.188 \\
\midrule
\multicolumn{7}{l}{\emph{Regresja}}\\
$D_J$ & -0.42 & 4/12 & 1.000 & +0.59 & 12/12 & \textbf{0.003} \\
$D_C$ & -0.27 & 4/12 & 1.000 & +0.69 & 12/12 & \textbf{0.003} \\
$D_Y$ & -0.46 & 4/12 & 1.000 & +0.58 & 12/12 & \textbf{0.003} \\
$U$ & +0.83 & 9/12 & 0.113 & +0.18 & 10/12 & 0.233 \\
$R_Y$ & -0.60 & 3/12 & 1.000 & +0.44 & 12/12 & \textbf{0.003} \\
$R_X$ & -0.65 & 3/12 & 1.000 & +0.29 & 12/12 & \textbf{0.003} \\
różnorodność predykcji & -0.28 & 5/12 & 1.000 & +0.77 & 12/12 & \textbf{0.003} \\
\bottomrule
\end{tabular}
\end{table}


\begin{table}[!htbp]
\centering\small
\caption{H2: korelacja Spearmana między średnią wartością miary a~wynikiem zespołu po 15 wariantach, liczona w~każdym zbiorze osobno (bagging, klasyfikacja). Podano medianę po zbiorach, liczbę zbiorów z~dodatnią korelacją i~wartość $p_H$ testu Wilcoxona z~korektą Holma (7 miar). Korelacja cząstkowa uwzględnia średnią jakość drzew.}
\label{tab:h2bag}
\begin{tabular}{lrrrrrr}
\toprule
& \multicolumn{3}{c}{\textbf{korelacja}} & \multicolumn{3}{c}{\textbf{cząstkowa}} \\
\cmidrule(lr){2-4}\cmidrule(lr){5-7}
\textbf{Miara} & $\tilde\rho$ & $\rho>0$ & $p_H$ & $\tilde\rho$ & $\rho>0$ & $p_H$ \\
\midrule
\multicolumn{7}{l}{\emph{Klasyfikacja}}\\
$D_J$ & -0.16 & 5/12 & 1.000 & +0.54 & 12/12 & \textbf{0.003} \\
$D_C$ & -0.09 & 5/12 & 1.000 & +0.60 & 12/12 & \textbf{0.003} \\
$D_Y$ & -0.24 & 5/12 & 1.000 & +0.54 & 12/12 & \textbf{0.003} \\
$U$ & +0.50 & 8/12 & 1.000 & +0.12 & 8/12 & 0.339 \\
$R_Y$ & -0.36 & 4/12 & 1.000 & +0.26 & 9/12 & 0.068 \\
$R_X$ & -0.37 & 5/12 & 1.000 & +0.41 & 10/12 & \textbf{0.028} \\
różnorodność predykcji & -0.16 & 5/12 & 1.000 & +0.62 & 12/12 & \textbf{0.003} \\
\bottomrule
\end{tabular}
\end{table}


\begin{figure}[!htbp]
\centering
\includegraphics[width=\textwidth]{06_h2_correlations.png}
\caption{H2 w~poszczególnych zbiorach (las losowy). Górny wiersz: korelacja Spearmana między średnią wartością miary a~jakością zespołu po 15 wariantach. Dolny wiersz: korelacja cząstkowa przy kontroli średniej jakości drzew.}
\label{fig:h2}
\end{figure}

% --------------------------------------------------------------------
\subsection{Miary próbek a~ranking wariantów}
\label{sec:wyniki-rangi}

Druga część PB3 dotyczy tego, czy miary obliczone przed uczeniem pozwalają wskazać lepszy wariant bez uczenia zespołów. Tabela~\ref{tab:miary-rangi} zestawia średnią wartość każdej miary ze średnią rangą wariantu. Najsilniej z~rangą wiąże się $U$ ($\rho=-0.90$ w~klasyfikacji i~$-0.70$ w~regresji): im więcej różnych obserwacji w~próbce, tym lepsza ranga. W~klasyfikacji istotne są także $D_Y$, $R_Y$, $R_X$ i~różnorodność predykcji; większa wartość tych miar oznacza gorszą rangę.

\begin{table}[!htbp]
\centering\small
\caption{Związek miar ze średnimi rangami wariantów: korelacja Spearmana między średnią wartością miary wariantu (po zbiorach) a~jego średnią rangą dla 15 wariantów. Korelacja ujemna oznacza, że wariant o~wyższej wartości miary ma lepszą rangę. Wartości $p_H$ podano po korekcie Holma.}
\label{tab:miary-rangi}
\begin{tabular}{lrrrr}
\toprule
& \multicolumn{2}{c}{\textbf{Klasyfikacja}} & \multicolumn{2}{c}{\textbf{Regresja}} \\
\cmidrule(lr){2-3}\cmidrule(lr){4-5}
\textbf{Miara} & $\rho$ & $p_H$ & $\rho$ & $p_H$ \\
\midrule
$D_J$ & +0.55 & 0.067 & +0.17 & 1.000 \\
$D_C$ & +0.45 & 0.089 & -0.01 & 1.000 \\
$D_Y$ & +0.64 & \textbf{0.031} & +0.21 & 1.000 \\
$U$ & -0.90 & \textbf{<0.001} & -0.70 & \textbf{0.027} \\
$R_Y$ & +0.79 & \textbf{0.003} & +0.42 & 0.592 \\
$R_X$ & +0.78 & \textbf{0.003} & +0.55 & 0.203 \\
różnorodność predykcji & +0.66 & \textbf{0.031} & +0.11 & 1.000 \\
\bottomrule
\end{tabular}
\end{table}


Rysunek~\ref{fig:miary-rangi} pokazuje jednak, że żadna pojedyncza miara nie porządkuje wariantów poprawnie. Przy tym samym $U=0.5$ cztery warianty S3 zajmują w~regresji rangi od 6 do 14, ponieważ $U$ nie uwzględnia pokrywania się próbek. Odległość Jaccarda je uwzględnia, ale jej związek z~jakością jest niemonotoniczny. Najgorsze warianty leżą na obu końcach skali. S3~$\gamma=0.75$ ma najmniejsze $D_J$, bo próbki są zbyt podobne, a~S1~$r=0.2$ i~S2~$s=0.18$ największe, bo próbki są różne, ale bardzo małe. To samo dotyczy różnorodności predykcji mierzonej po uczeniu. Najlepsze warianty mają duże $U$ i, przy tym $U$, możliwie duże $D_J$. W~tym obszarze leży wariant odniesienia.

\begin{figure}[!htbp]
\centering
\includegraphics[width=\textwidth]{10_measures_vs_ranks.png}
\caption{Średnia wartość miary wariantu (po zbiorach) a~jego średnia ranga w~teście Friedmana. Kolumny: udział różnych obserwacji $U$, odległość Jaccarda $D_J$ oraz, dla porównania, różnorodność predykcji mierzona po uczeniu.}
\label{fig:miary-rangi}
\end{figure}

% --------------------------------------------------------------------
\subsection{Eksperyment dodatkowy: wpływ konkretnego losowania}
\label{sec:wyniki-losowania}

Przy ustalonym wariancie (bagging, wariant odniesienia) 20 zespołów różniących się tylko ziarnem losowania próbek nie wykazuje związku między miarami próbek a~jakością zespołu (Tabela~\ref{tab:losowania}, Rysunek~\ref{fig:losowania}). Mediany korelacji mieszczą się w~przedziale $\pm0.1$, znak zmienia się między zbiorami i~żaden test nie jest istotny. Na podstawie miar próbek nie da się więc wybrać lepszego z~20 wylosowanych zespołów.

\begin{table}[!htbp]
\centering\small
\caption{Eksperyment dodatkowy: korelacja Spearmana między miarą próbek a~jakością zespołu dla 20 losowań próbek w~tym samym podziale (bagging, wariant S1~$r=1$). Podano medianę po zbiorach z~median po 5 podziałach oraz wartość $p_H$ testu Wilcoxona z~korektą Holma (6 miar).}
\label{tab:losowania}
\begin{tabular}{lrrrr}
\toprule
& \multicolumn{2}{c}{\textbf{Klasyfikacja}} & \multicolumn{2}{c}{\textbf{Regresja}} \\
\cmidrule(lr){2-3}\cmidrule(lr){4-5}
\textbf{Miara} & $\tilde\rho$ & $p_H$ & $\tilde\rho$ & $p_H$ \\
\midrule
$D_J$ & -0.043 & 1.000 & +0.044 & 1.000 \\
$D_C$ & -0.026 & 1.000 & -0.011 & 1.000 \\
$D_Y$ & -0.046 & 0.463 & +0.040 & 1.000 \\
$U$ & +0.023 & 1.000 & +0.024 & 1.000 \\
$R_Y$ & -0.092 & 1.000 & +0.030 & 1.000 \\
$R_X$ & -0.057 & 0.647 & -0.018 & 1.000 \\
\bottomrule
\end{tabular}
\end{table}


\begin{figure}[!htbp]
\centering
\includegraphics[width=\textwidth]{11_draws.png}
\caption{Eksperyment dodatkowy: w~każdym zbiorze mediana (po 5 podziałach) korelacji Spearmana między średnią miarą próbek a~jakością zespołu dla 20 losowań próbek w~baggingu (wariant odniesienia).}
\label{fig:losowania}
\end{figure}

% --------------------------------------------------------------------
\subsection{Las losowy a~bagging}
\label{sec:rf-bagging}

W~klasyfikacji powtórzyliśmy eksperyment dla baggingu, czyli drzew rozważających w~każdym węźle wszystkie cechy (Sekcja~\ref{sec:model}). Także tu test Friedmana odrzuca równoważność wariantów ($\chi^2_F=55.9$, $F_{ID}=5.49$, $p<0.001$), a~żaden wariant nie jest istotnie lepszy od wariantu odniesienia. Rysunek~\ref{fig:rf-bagging} pokazuje dwie różnice między konfiguracjami.

Bagging jest znacznie bardziej wrażliwy na pokrywanie się próbek. S3~$\gamma=0.75$ obniża macro-F1 o~medianę 1.48\pp{}, a~w~lesie losowym o~0.48\pp{}. S3~$\gamma=0.5$ i~$\gamma=0.75$ przegrywają z~wariantem odniesienia we wszystkich 12 zbiorach ($p_H=0.007$). Korelacje cząstkowe z~H2 są w~baggingu istotne we wszystkich zbiorach (Tabela~\ref{tab:h2bag}), a~związek badany w~H1 jest silniejszy ($\tilde\rho=0.033$ zamiast 0.018, 12 z~12 zbiorów, Tabela~\ref{tab:h1}).

Las losowy jest natomiast bardziej wrażliwy na małe próbki. S1~$r=0.2$ obniża w~nim macro-F1 o~1.62\pp{}, a~w~baggingu o~1.16\pp{}. W~baggingu umiarkowanie mniejsze próbki (S1~$r=0.6$, S2~$s=0.45$) mają nawet minimalnie dodatnią medianę, choć mieszczącą się w~ROPE.

Predykcje drzew lasu losowego są bardziej różne niż predykcje drzew baggingu we wszystkich 12 zbiorach (prawy panel Rysunku~\ref{fig:rf-bagging}). Losowanie cech jest więc drugim, niezależnym od losowania obserwacji źródłem różnorodności. Częściowo zastępuje ono różnorodność próbek, ale osłabia drzewa, dlatego dodatkowe zmniejszenie próbki szkodzi lasowi losowemu bardziej niż baggingowi.

\begin{figure}[!htbp]
\centering
\includegraphics[width=\textwidth]{07_rf_vs_bagging.png}
\caption{Las losowy a~bagging (klasyfikacja). Lewy panel: mediana po zbiorach różnicy macro-F1 każdego wariantu względem wariantu odniesienia w~tej samej konfiguracji. Prawy panel: średnia różnorodność predykcji drzew (niezgodność) w~wariancie odniesienia; linia przerywana oznacza równość.}
\label{fig:rf-bagging}
\end{figure}
